\section{Introduction}

Reinforcement learning (RL) which is increasingly used to train very large machine learning models has two training paradigms: online and offline. While in online RL a policy model is trained while interacting with the environment, offline RL trains a model on data points collected from multiple trajectories of interactions with the environment and external entities (e.g. humans, AI agents). In many applications % examples
online RL is either not possible, very expensive or not scalable due the requirement of a dedicated agent interacting with the environment for the duration of model training. A key benefit of offline RL is its ability to ingest large amounts of diverse training data~\citep{fu2021d4rldatasetsdeepdatadriven}, and consequently, offline RL has become popular for training large models for tasks in for e.g. natural language processing, computer vision, robotics etc~\citep{levine2020offlinereinforcementlearningtutorial}. 


However, the use of large scale datasets presents challenges related to their storage, management as well the computational expense incurred in their use for training multiple models with different hyperparameters and architectures.  
One way to mitigate this issue is to create a smaller synthetic dataset derived from the training dataset i.e., a \emph{distillation} of the latter. Dataset distillation (DD) for supervised learning has been extensively studied in previous works (e.g. \cite{Wangetal18,sachdeva2023data,10273632}) which have developed a variety of methods based on matching the loss gradients or feature-embeddings between the training and synthetic datasets by optimizing the latter. This is done either with respect to networks being trained simultaneously using bi-level optimization, or a fixed collection of sampled networks. 
Distinct from sampling, coreset or random projection methods, DD techniques generate synthetic data via optimization with respect to networks which may also be trained in the process. These DD techniques have been shown to perform well on real-world supervised datasets in terms of the quality as well as the size of the generated syntheic dataset as the latter is explicitly optimized.  

While some works (e.g. \cite{nguyen2021dataset,chen2024provable}) have also provided theoretical performance guarantees, most DD techniques are heuristics albeit with empirically observed gains.  


For offline RL, a few recent works have proposed methods using behavioral cloning (BC)  for distilling an offline training dataset comprising trajectories of (state, action)-pairs. In \cite{lei2024offline}, the authors propose optimizing the synthetic dataset using policy-based BC loss leveraging action-value weights learnt from offline RL. On the other hand, the method proposed by \cite{light2024datasetdistillationofflinereinforcement} extends the matching loss gradients approach to offline RL. In particular, it optimizes a synthetic dataset to match the BC loss gradients on the offline and synthetic data, where the gradients are with respect to the parameters of sampled action predictor networks. 


As can be seen, the state of research into offline RL dataset distillation, while nascent, is also unsatisfying. Firstly, the BC does not leverage the observed reward that is usually available in offline RL datasets, and typically only perform well when the training dataset is generated by expert policies. Secondly, there is a sparsity of theoretical performance guarantees for the DD techniques for offline RL as well as supervised learning DD. For linear regression, the work of \cite{chen2024provable} proves efficiency guarantees assuming however that a trained model is available, while the work of \cite{nguyen2021dataset} -- while not training a model -- proves convergence only of the synthetic feature-vectors for given synthetic labels. For offline RL, the work of \cite{lei2024offline} provides analytical guarantees, while requiring an trained policy model on the given training dataset. To the best of our knowledge, there are no provably efficient DD algorithms known without model training for either supervised learning or offline RL. 



{\bf \underline{Our Contributions.}}  In this work we make progress towards bridging the above gaps in our understanding through the following contributions:

{\bf Dataset Distillation in supervised regression.} We propose an algorithm minimizing the squared difference of MSE losses between the training and the synthetic datasets with respect to a fixed set of randomly sampled models.  While our algorithm is along the line of previous DD approaches using loss, gradient loss or embedding matching techniques, our contribution is to prove that it  admits  efficiency and performance guarantees for linear regression. In particular, we prove that the optimization objective is convex and tractable, and in $d$-dimensions optimizing the synthetic dataset w.r.t. $\tilde{O}(d^2)$\footnote{$\tilde{O}$ hides polylogarithmic factors} randomly sampled regressors suffices to guarantee that any bounded linear regressor has approximately the same MSE loss on the training and synthetic datasets. Therefore, the optimum linear regressor on the synthetic dataset is close to that on the training dataset. This is the first performance and efficiency guarantee for supervised DD without model training. \\ 
We also prove a lower bound of $\Omega(d^2)$ on the number of randomly sampled linear regressors to obtain a good quality synthetic dataset. This result shows the tightness of our algorithmic guarantees.


{\bf Dataset Distillation in offline RL.} We extend the above approach for dataset distillation in supervised regression to offline RL DD by matching the Bellman loss on the training and synthetic datasets  w.r.t. a collection of randomly sampled linear $Q$-value predictors over $\R^d$. 
The performance guarantees are in a similar vein as the supervised regression case, though more involved due to the $\max$ term in the Bellman Loss formulation. 

Nevertheless, we obtain a performance guarantee showing that minimizing to a small value the difference of the Bellman Losses with respect to $\tn{exp}(O(d\log d))$ randomly sampled $Q$-value predictors (without model training) generates a synthetic dataset whose Bellman Loss w.r.t. any linear bounded $Q$-value predictor is a close approximation to that on the training dataset. Consequently, the Bellman Loss minimizer $Q$-value predictor on the synthetic dataset is close to optimal on the training dataset. 

{\bf Offline RL with decomposable feature-maps.} 
We consider a natural setting where the state-action feature-map is the sum of  individual feature-maps of the state and action. In this scenario, we prove that our method requires only $\tilde{O}(d^2)$ sampled linear $Q$-value predictions, instead of the exponentially many required in the general case. Additionally, we show that the finding the synthetic dataset is a tractable convex optimization when the state and action feature-maps are linear.



{\bf Empirical Evaluations.} We conduct experiments on supervised regression and offline RL datasets to validate our proposed algorithms, showing that our techniques obtain improved performance over standard baselines, both in terms of quality and size of the generated datasets. Notably, our experiments demonstrate the effectiveness of our algorithms using very few sampled models in both supervised regression and offline RL settings. Although the guarantees are proven for linear models, we show that our algorithms work well in practice with non-linear neural networks.





